% 05_run.tex - run the coupled model, read the summary, units, balance check
\section{Run the coupled model}

\step{Run Raven}\label{st:runa}
\term{t4_run_a.txt}
\see \msg{Exiting Gracefully: Successful Simulation}, after about \V{t4:runsec}~seconds. The files as they are in this step are in
\cmd{\$WS/model/step\_07}. Now read the summary:
\term{t4_summary.txt}
\begin{tocheck}
\begin{itemize}
\item \cmd{existing model} names the model, its grid and its units (\emph{feet, days}).
\item \cmd{packages taken over} and \cmd{packages in balance}. The packages in balance are the model's own packages that
  stay (here the well and the head boundary). Their net flow is the column \cmd{other model packages} of
  \cmd{GWBudget.csv}.
\item \cmd{HRU coverage}: \V{t4:coverage}\% of the coupled HRUs' area lies on active cells. The rest lies outside the
  model, whose edges do not follow the HRU outlines and cut off the outer parts of the slope HRUs. A low value usually
  means a wrong projection or wrong length units.
\item The second part of the file lists each HRU with \cmd{linked\_fraction\_of\_polygon}. The slope HRUs 3, 7 and 8 are
  linked at \V{t4:link3}, \V{t4:link7} and \V{t4:link8}\%, the others at \V{t4:linkothers}\% or more.
\end{itemize}
\end{tocheck}
\begin{trouble}
\begin{itemize}
\item \emph{``\dots would add recharge or evapotranspiration beside Raven's''}: name every recharge, ET and UZF package
  in \cmd{:RavenTakesOver} or \cmd{:KeepPackages}.
\item \emph{``stream package entry with subbasin ID \dots{} - no such subbasin''}: the \cmd{SUBBASIN} values must be
  subbasin numbers of the \cmd{.rvh}.
\item \emph{``\dots is not a whole number of Raven time steps''} or \emph{``the MODFLOW 6 model ends before Raven's
  run''}: check \cmd{:MF6StartDate}, and \cmd{:StartDate} and \cmd{:Duration} in the \cmd{.rvi}.
\item Commands for a grid that Raven builds (\cmd{:GridCellSize}, \cmd{:LandSurface}, \cmd{:Wells} and others) stop the
  run here: the model has its own grid and wells.
\end{itemize}
\end{trouble}

\step{See what Raven changed in the working copy}
The working copy in \cmd{out\_coupled/mf6/} shows what Raven did to the model's name file:
\term{t4_mf6nam.txt}
\why the packages that Raven took over are commented out, and Raven's own packages are added at the end, each marked as
added by Raven.

\paragraph{Units.} The model works in feet and days. Raven converts at the interface (1~ft = 0.3048~m,
1~ft\textsuperscript{3} = 0.0283168~m\textsuperscript{3}). It writes the recharge, seepage, river and package flows to
\cmd{GWBudget.csv} in m\textsuperscript{3}/d and the heads in metres, as for a model that Raven builds. A model in seconds
or in centimetres is converted in the same way. A model that does not give its units is taken as metres and days.

\step{Check the water balance}
\term{t4_audit_a.txt}
\begin{tocheck}
The recharge line is 0 (all the water Raven sent reached the model's cells, in feet), the cumulative residual is
\V{t4:a:auditRes}, and there are \V{t4:a:auditUnconv} unconverged steps.
\end{tocheck}
\why the mean flows, in m\textsuperscript{3}/d, are: recharge \V{t4:a:rech}, seepage to Raven \V{t4:a:seep}, discharge to
the creek \V{t4:a:toriv}, leakage from the creek \V{t4:a:fromriv}, the model's own well and boundary \V{t4:a:other}, and
storage release \V{t4:a:stor}.
\begin{note}[Look at the storage]
The storage release is large and positive: \V{t4:a:stor1}~m\textsuperscript{3}/d on average in the first year, and still
\V{t4:a:stor3} in the third. The model's first period is a steady state. It is solved with the recharge of Raven's first
day, \V{t4:a:rechday1}~m\textsuperscript{3}/d (\V{t4:a:rechday1mm}~mm over the coupled HRUs), which is high because the
initial soil moisture of the \cmd{.rvc} lets the subsoil drain fast. So the model starts with too much water and drains for
years. In a real study, start Raven from soil states after a spin-up (the \cmd{solution.rvc} of an earlier run), or let the
model's steady period see a typical recharge.
\end{note}
